% !TEX program = pdflatex
% T-DMB end-to-end open-source receiver — arXiv preprint
%
% Overleaf:
%   1. Create new project, upload this file + paper/figures/*
%   2. Compile with pdflatex (preferred) or xelatex
%   3. The bibliography uses biber + biblatex — Overleaf builds it automatically
%
% Recommended template: arXiv-style (single column) for the preprint;
% switch to IEEEtran two-column for a journal/conference resubmission.

\documentclass[11pt,a4paper]{article}

\usepackage[a4paper,margin=1in]{geometry}
\usepackage[utf8]{inputenc}
\usepackage[T1]{fontenc}
\usepackage{lmodern}
\usepackage{microtype}
\usepackage{graphicx}
\usepackage{booktabs}
\usepackage{array}
\usepackage{caption}
\usepackage{subcaption}
\usepackage{xcolor}
\usepackage{listings}
\usepackage{tikz}
\usepackage{hyperref}
\usepackage{url}
\usepackage[backend=biber,style=ieee,sorting=none]{biblatex}
\usepackage{amsmath}
\usepackage{enumitem}

\hypersetup{
    colorlinks=true,
    linkcolor=blue!60!black,
    urlcolor=blue!50!black,
    citecolor=blue!50!black
}

\lstdefinestyle{code}{
    basicstyle=\ttfamily\footnotesize,
    breaklines=true,
    columns=fullflexible,
    keepspaces=true,
    frame=single,
    rulecolor=\color{black!30},
    backgroundcolor=\color{black!3},
    numbers=none
}
\lstset{style=code}

\addbibresource{refs.bib}

% --------------------------------------------------------------------
\title{An End-to-End Open-Source Software Receiver for Korean
       Terrestrial DMB on Low-Cost SDR Hardware}

\author{%
  Seonggeun Yoo\textsuperscript{1}\thanks{Correspondence:
    \texttt{orcogre@vendit.co.kr}} \quad
  Claude\textsuperscript{2}\thanks{Claude is a large language model
    developed by Anthropic.  It contributed to this work as a
    paired-programming and writing assistant: reverse-engineering
    the Korean T-DMB SL-packet header structure, drafting the
    Python implementation of the sync-aligned outer FEC, locating
    the embedded AVC \texttt{DecoderConfigurationRecord} in the
    BIFS stream, and producing the initial draft of this manuscript.
    All design decisions, experimental choices, and final editorial
    control rest with the primary author.}\\
  \small\textsuperscript{1}Independent Researcher, Seoul,
                            Republic of Korea\\
  \small\textsuperscript{2}Anthropic, San Francisco, USA\\
}

\date{\today}

% --------------------------------------------------------------------
\begin{document}
\maketitle

\begin{abstract}
Terrestrial Digital Multimedia Broadcasting (T-DMB) has been in continuous
operation in the Republic of Korea since December 2005, yet two decades on
there exists no \emph{complete} open-source software stack that takes a
T-DMB RF signal from a low-cost software-defined radio all the way to a
decoded video frame.  Prior efforts fall into two camps: (i) DAB receivers
such as \texttt{welle.io} and \texttt{eti-cmdline} that stop after producing
Ensemble Transport Interface (ETI) frames, and (ii) the recent GNU Radio +
patched FFmpeg pipeline of Eo and Bahk~\cite{eo2024tdmb} which decodes
audio/video but is deeply tied to GNU Radio's graphical companion and a set
of out-of-tree patches.  This paper presents a third path: a self-contained
Python package implementing the full T-DMB receive chain --- ETI parsing,
Fast Information Channel (FIC) decoding, Reed-Solomon (RS) plus
convolutional time-deinterleaver outer FEC, MPEG-4 Sync Layer (SL) packet
demultiplexing, and H.264 elementary-stream reconstruction --- with the
final video decoding stage delegated to an unmodified build of the
\texttt{dmb-oss} FFmpeg fork.  We verify the chain on a 30\,MB
on-air capture of YTN DMB (channel K8B, 183.008\,MHz) using an
Airspy~Mini~\cite{airspy} and a consumer indoor antenna: 87.3\% of
Reed-Solomon blocks decode cleanly, every SL packet header parses, and
ffmpeg renders QVGA H.264 Baseline video frames containing recognisable
Korean broadcast content.  We identify and fix a previously unreported
synchronisation bug that prevented the conventional ``deinterleave-then-find-sync''
ordering from converging on real captures, and we document an empirical
discovery that the AVC \texttt{DecoderConfigurationRecord} in Korean
T-DMB broadcasts is embedded inside the Scene-Description / BIFS stream
rather than the Object-Descriptor stream as MPEG-4 Systems would suggest.
The code, captures, and per-capture metadata sidecars are released on
GitHub.
\end{abstract}

\noindent\textbf{Keywords:}
Software-defined radio, T-DMB, DAB, MPEG-4 SL packets, Reed-Solomon,
convolutional interleaver, open source.

% ====================================================================
\section{Introduction}\label{sec:intro}

Korean Terrestrial DMB extends the European Eureka-147 DAB physical
layer~\cite{etsi300401} with an additional Reed-Solomon outer code and a
Forney convolutional time interleaver, in order to deliver MPEG-4
AVC~/~H.264 video and MPEG-4 BSAC audio over MPEG-2 transport streams
multiplexed into DAB sub-channels~\cite{ttak0026,etsi102427,etsi102428}.
The service has been in commercial operation since December 2005, was
deployed nationwide, and remains in service in 2026 despite the rise of
smartphone streaming.  Its specifications are publicly available, the
modulation is the same DAB Mode~I OFDM used across Europe, and dedicated
DAB demodulators exist as open source (notably \texttt{welle.io}, the
\texttt{gr-dab} GNU Radio out-of-tree module, and Jan van Katwijk's
\texttt{eti-stuff}).  In spite of this, until very recently no published
software released under a free licence has been able to take a T-DMB RF
capture and produce a playable video file end-to-end without proprietary
tools.

The closest precedent is the recent ICTC paper by
Eo and Bahk~\cite{eo2024tdmb} (with their \texttt{dmb-oss} repository),
which demonstrated playback by chaining \texttt{gr-dab} with a patched
FFmpeg.  Their work is the canonical reference for the Korean T-DMB
encoding chain and we rely on it heavily.  However the released artefact
is a GNU~Radio flowgraph plus a set of FFmpeg patches; it requires the
graphical GNU~Radio Companion, has not been packaged for headless or
library use, and reports several practical gotchas (alignment of the
deinterleaver, missing composition timestamps) without
all-in-one tooling to mitigate them.

This paper describes an independent re-implementation with three goals:
\begin{enumerate}[leftmargin=*,nosep]
\item \textbf{Library-first.}  Every stage is a Python module
       (\texttt{tdmb.eti}, \texttt{tdmb.fec}, \texttt{tdmb.msc}, ...) with a
       documented API, so individual layers can be unit-tested and reused
       (e.g.\ a pure-Python Forney deinterleaver that decodes the same
       reference capture as our pipeline in under a minute).
\item \textbf{Reproducible captures.}  Every RF capture writes a JSON
       sidecar recording the channel, frequency, sample rate, gain stages,
       bias-tee state, antenna model, and capture location.  We retro-fit
       sidecars for our reference captures so that the artefact bundle is
       self-describing.
\item \textbf{One-shot end-to-end driver.}  A single command
       (\texttt{python tools/decode\_video.py <capture>.eti}) traverses
       the full chain and emits an MP4 plus IDR thumbnail PNGs.  No GUI
       and no FFmpeg patches at run time are required (we link against a
       \emph{pre-built} unmodified \texttt{dmb-oss} FFmpeg for the final
       H.264~/~BSAC decode).
\end{enumerate}

\subsection{Contributions}
\begin{itemize}[leftmargin=*,nosep]
\item A pure-Python implementation of the T-DMB outer FEC chain
      (RS(204,188) over $GF(2^8)$ with DVB-T polynomial plus a
      $12\times17$ Forney deinterleaver) that achieves 87.3\% block
      success on a reference Korean K8B capture with FIB CRC 100\%.
\item Identification of, and a fix for, a subtle phase-alignment issue:
      feeding the deinterleaver from an arbitrary offset of the MSC
      byte stream silently scrambles every 204-byte block even though
      the TS sync byte~(\texttt{0x47}) remains at the correct
      cadence.  Pre-aligning to the empirically-determined sync byte
      offset (160 within the 204-byte cycle on our captures) restores
      bit-exact deinterleaving.
\item A reverse-engineered SL packet header layout for the Korean
      T-DMB video stream (a fixed 9-byte header per Access Unit:
      6-bit flags, 33-bit DTS, 33-bit CTS, leading flag byte
      \texttt{0xCF}) verified against \(>\,3000\) PES packets.
\item Empirical observation that the AVC \texttt{DecoderConfigurationRecord}
      for the video elementary stream is carried in the Scene-Description
      (table\_id~\texttt{0x05}) section stream rather than the
      Object-Descriptor (table\_id~\texttt{0x04}) stream used by the
      reference MPEG-4 Systems specification~\cite{iso14496-1}.  We
      provide a heuristic locator and confirm it on the YTN DMB
      ensemble.
\item Release of an Overleaf-ready paper source, a Python package, a
      30\,MB reference capture with its metadata sidecar, and twenty
      decoded IDR thumbnails from said capture.
\end{itemize}

% ====================================================================
\section{System Overview}\label{sec:system}

\begin{figure}[t]
\centering
% Pipeline diagram for the T-DMB end-to-end receiver paper.
% Drop-in TikZ block: usepackage{tikz} in main.tex.
\usetikzlibrary{positioning,arrows.meta,shapes.geometric,calc}

\begin{tikzpicture}[
    >=Stealth,
    node distance=4mm and 4mm,
    every node/.style={font=\footnotesize},
    stage/.style={draw, rounded corners=2pt, minimum height=8mm,
                  minimum width=27mm, align=center, fill=blue!4},
    py/.style   ={draw, rounded corners=2pt, minimum height=8mm,
                  minimum width=27mm, align=center, fill=green!8},
    ext/.style  ={draw, rounded corners=2pt, minimum height=8mm,
                  minimum width=27mm, align=center, fill=orange!10},
    io/.style   ={draw, dashed, rounded corners=2pt, minimum height=7mm,
                  minimum width=27mm, align=center, fill=black!3,
                  font=\footnotesize\ttfamily},
]

\node[io]                          (rf)    {RF capture\\(Airspy Mini)};
\node[stage, right=of rf]          (eti)   {eti-cmdline\\(OFDM + Viterbi)};
\node[py, right=of eti]            (fic)   {tdmb.eti.fic\\(FIB $\to$ ensemble)};
\node[py, below=of eti]            (msc)   {tdmb.msc\\(extract sub-channel)};
\node[py, below=of msc]            (outer) {tdmb.fec.outer\\\textbf{sync-aligned}\\RS(204,188)+TI};
\node[py, right=of outer]          (sl)    {SL packet demux\\(strip 9-byte header)};
\node[py, right=of sl]             (avcc)  {AVCC extractor\\(scan SD/BIFS)};
\node[io, right=of avcc]           (h264)  {H.264 Annex B\\elementary stream};
\node[ext, below=of h264]          (ff)    {dmb-oss FFmpeg\\(unmodified)};
\node[io, left=of ff]              (mp4)   {MP4 / PNG\\thumbnails};

\draw[->] (rf)    -- (eti);
\draw[->] (eti)   -- (fic);
\draw[->] (eti)   -- (msc);
\draw[->] (msc)   -- (outer);
\draw[->] (outer) -- (sl);
\draw[->] (sl)    -- (avcc);
\draw[->] (avcc)  -- (h264);
\draw[->] (h264)  -- (ff);
\draw[->] (ff)    -- (mp4);

\node[font=\scriptsize, anchor=west, text=black!60] at ($(rf.south)+(-2mm,-2mm)$)
      {\textit{Stage 1--2 (delegated)}};
\node[font=\scriptsize, anchor=west, text=black!60] at ($(outer.south)+(-2mm,-2mm)$)
      {\textit{Stage 3--6 (this work, Python)}};
\node[font=\scriptsize, anchor=west, text=black!60] at ($(ff.south)+(-2mm,-2mm)$)
      {\textit{Stage 7 (delegated)}};

\end{tikzpicture}

\caption{End-to-end pipeline from RF to decoded H.264 frame.  Stages 1
and 2 (DAB demodulation through ETI) are delegated to a patched build
of \texttt{eti-cmdline-airspy} (one-line patches to the band-handler
table to recognise Korean K-prefix channels and to the Airspy front-end
to accept environment-variable gain overrides).  Stages 3--7
(ETI~$\rightarrow$~FIC,
MSC,~RS+TI,~SL,~AVCC,~H.264~ES) are pure Python.  The final video
decode is performed by \texttt{dmb-oss} FFmpeg with no run-time
modifications.}
\label{fig:pipeline}
\end{figure}

Figure~\ref{fig:pipeline} summarises the chain.  We take captures with
an Airspy~Mini SDR at 3\,MSPS or 6\,MSPS, tuned to one of the Korean
K-block centre frequencies listed in Table~\ref{tab:k-block}.
The 12-channel K-block raster spaces VHF Band~III differently from the
ETSI standard~\cite{etsi300401} (1.728\,MHz steps obtained by dividing
each 6\,MHz analogue TV channel into three equal sub-channels A/B/C),
so the band-handler table inside \texttt{eti-cmdline} must be patched
to add the \texttt{K}-prefixed entries before any Korean capture
will tune correctly.

\begin{table}[t]
\centering
\caption{Korean K-block channel raster and known operators for the
Seoul metropolitan area.}
\label{tab:k-block}
\begin{tabular}{llll}
\toprule
Channel & Centre frequency & Spacing & Operator (Seoul) \\
\midrule
K8B  & 183.008\,MHz & 1.728\,MHz & YTN DMB \\
K12A & 205.280\,MHz & 1.728\,MHz & MBC DMB \\
K12B & 207.008\,MHz & 1.728\,MHz & U-KBS \\
K12C & 208.736\,MHz & 1.728\,MHz & SBS u \\
\bottomrule
\end{tabular}
\end{table}

% ====================================================================
\section{Outer FEC: Sync-Aligned Deinterleaving}\label{sec:rs}

T-DMB augments the DAB physical layer with an outer protection step
specified in ETSI TS~102\,427~\cite{etsi102427}: each 188-byte MPEG-2
transport-stream packet is first Reed-Solomon encoded into a
204-byte codeword (RS(204,188) shortened from RS(255,239), with
$\alpha=2$, primitive polynomial $0x11d$, FCR$=0$ -- identical to
DVB-T), and the resulting codeword stream is then byte-wise
convolutional-interleaved using a Forney scheme with $N=12$ branches
and depth $M=17$.

At the receiver, after Viterbi inner-FEC decoding and energy
descrambling we obtain a stream of MSC sub-channel bytes.  A naive
implementation runs the Forney deinterleaver immediately, then scans
the deinterleaved output for the 204-byte alignment of the TS sync
byte~(\texttt{0x47}).  On our reference K8B capture this approach
recovered \emph{zero} successful Reed-Solomon blocks out of 25\,953
attempts.

The cause is that the conv-interleaver branches are indexed modulo
$N=12$ based on stream position.  Byte~0 of every transmit-side RS
block goes through branch~0 (zero delay) and so emerges at transmit
positions $0, 204, 408, \ldots$ --- but the very first MSC byte we
\emph{observe} in our capture is not byte~0 of any RS block; it is
some arbitrary phase $\phi$ within the 204-byte cycle.  If we feed that
byte into the receive-side deinterleaver as ``branch~0,'' the branch
indices are off by $\phi \bmod 12$ for the entire stream, and the
deinterleaver permutes every byte to the wrong slot.  Crucially the
0x47 sync byte still emerges at the same cadence (because every
$204$-byte block also has its first byte on branch~0 of the receive
deinterleaver, by coincidence of $204=N\cdot M$), so a sync-scan
\emph{after} deinterleaving cannot distinguish the misaligned case
from the aligned one.  This matches an observation in passing by
Eo and Bahk~\cite{eo2024tdmb}: ``the deinterleaver requires the sync
byte of a TS packet to be the first in an input byte chunk.''

Our fix is to perform sync detection \emph{before} the deinterleaver.
We accumulate $\sim$50 RS blocks ($\sim$10\,kB) of raw MSC bytes,
compute the count of \texttt{0x47} occurrences at each of 204 candidate
phases, and require the best phase to be (a)~hit by $\geq 40\%$ of
blocks and (b)~$\geq 5\times$ the second-best phase.  Once locked, we
discard the bytes preceding the chosen phase and feed the remainder
into a freshly-constructed Forney deinterleaver.  After the
$(N{-}1)\cdot M\cdot N = 2244$-byte settling transient the
deinterleaver emits 204-byte RS codewords ready for systematic
correction.

Table~\ref{tab:rs} reports the result on the YTN-DMB K8B reference
capture (FIB CRC 100\%, $\sim$2\,minutes of content).

\begin{table}[t]
\centering
\caption{Outer-FEC decode statistics on the K8B reference capture,
sub-channel~1 (mYTN video).  ``Sync-aligned'' is our fix; the legacy
column is the implementation that runs the deinterleaver before sync
search.}
\label{tab:rs}
\begin{tabular}{lrr}
\toprule
Metric & Legacy & Sync-aligned (this work) \\
\midrule
RS blocks total                 & 25\,953 & 25\,953 \\
RS ok (zero errors)             & 0       & 0 \\
RS corrected ($\leq8$ bytes)    & 0       & 22\,652 \\
RS uncorrectable                & 25\,953 & 3\,301 \\
\textbf{Success rate}           & \textbf{0.0\,\%} & \textbf{87.3\,\%} \\
\bottomrule
\end{tabular}
\end{table}

The remaining 13\% of uncorrectable blocks correspond to bursty
errors that exceed the $t=8$ correction limit; these are concentrated
near CIF boundaries and during occasional fade-induced SNR dips.
At the slice level, FFmpeg's macroblock-level error concealment
successfully reconstructs visually-coherent frames.

% ====================================================================
\section{SL Demultiplexing}\label{sec:sl}

The transport stream emerging from the RS decoder carries four logical
PIDs (Table~\ref{tab:pids}).  PID~\texttt{0x113} contains the video
elementary stream packetised first into MPEG-4 SL packets and then
into MPEG-2 PES with \texttt{stream\_id} = \texttt{0xFA}
(``MPEG-4 SL-packetised stream'').  Each PES carries exactly one SL
packet, and each SL packet contains exactly one H.264 NAL unit
(without start-code prefix).

\begin{table}[t]
\centering
\caption{PID assignment recovered from the YTN DMB K8B ensemble.}
\label{tab:pids}
\begin{tabular}{lll}
\toprule
PID & Carried stream & Packet count in 2 min \\
\midrule
\texttt{0x100} & PMT                              & 212 \\
\texttt{0x111} & ObjectDescriptor sections (\texttt{tid}=0x04) & 224 \\
\texttt{0x112} & SceneDescription sections (\texttt{tid}=0x05) & 220 \\
\texttt{0x113} & Video PES, SL-packetised         & 16\,397 \\
\texttt{0x114} & Audio PES, SL-packetised (BSAC)  & 4\,916 \\
\texttt{0x1FFF} & Null padding                     & 457 \\
\bottomrule
\end{tabular}
\end{table}

\subsection{SL header layout}

Reading the ISO/IEC 14496-1 specification of
\texttt{SLPacketHeader()} \cite{iso14496-1} along with the FFmpeg
\texttt{read\_sl\_header()} implementation in \texttt{dmb-oss}, the
maximum-length AU-start header for the Korean T-DMB profile contains:

\begin{align*}
\text{header bits} \;=\; & 1 \;(\text{au\_start}) + 1 \;(\text{au\_end}) + 1 \;(\text{padding}) + \\
                         & 1 \;(\text{rand\_acc\_pt}) + 1 \;(\text{dts\_flag}) + 1 \;(\text{cts\_flag}) + \\
                         & 33 \;(\text{DTS}) + 33 \;(\text{CTS}) \\
                       =\; & 72 \; \text{bits} \;=\; 9 \; \text{bytes exactly}.
\end{align*}

We confirm this empirically on \emph{every one} of the 3\,135 SL
packets in our reference capture: each PES body starts with the byte
\texttt{0xCF} (=$0b11001111$, decoding to
\{au\_start,au\_end\}=\{1,1\}, \{padding,rand\_acc\_pt\}=\{0,0\},
\{dts\_flag,cts\_flag\}=\{1,1\}, plus two high bits of DTS), followed
by 8 more bytes of timestamp fields, after which byte~9 is the NAL
header of a single H.264 NAL unit.

\subsection{NAL stream}

After stripping the 9-byte SL header from every PES we obtain a stream
of bare H.264 NAL bytes.  The NAL header histogram on our reference
capture is dominated by non-IDR slices (3\,038) with periodic IDR
slices (97), giving a Group-of-Pictures (GoP) length of roughly~32 ---
typical for a 30\,fps QVGA broadcast.  No SPS or PPS NAL units are
present in the elementary stream, which is consistent with
AVCC-style storage where decoder configuration is carried out-of-band.

% ====================================================================
\section{Locating the AVC Decoder Configuration Record}
\label{sec:avcc}

In standard MPEG-4 Systems, the H.264 \texttt{AVCDecoderConfigurationRecord}
(which contains the SPS and PPS NAL units required to bootstrap the
decoder) is transported inside the Object-Descriptor stream as part of a
\texttt{DecoderSpecificInfo} descriptor~\cite{iso14496-1}.  In the YTN
DMB capture we observe the opposite: the OD-stream sections on
PID~\texttt{0x111} contain only DMB-private commands with user-defined
tags (\texttt{0xC0}--\texttt{0xFE}, the ``Annex K user-private'' range);
the actual AVCC bytes appear inside the Scene-Description /
BIFS section stream on PID~\texttt{0x112}.

Our locator scans the SD-stream byte stream for the AVCC signature
$$\texttt{01}\;\langle\textit{profile}\rangle\;\langle\textit{compat}\rangle
   \;\langle\textit{level}\rangle\;\texttt{FF}\;\texttt{E1}\;
   \langle\textit{SPS\_len}\rangle\;\langle\textit{SPS}\rangle\;
   \texttt{01}\;\langle\textit{PPS\_len}\rangle\;\langle\textit{PPS}\rangle$$
constraining profile to a known H.264 value
(\texttt{0x42}/\texttt{0x4D}/\texttt{0x58}/\texttt{0x64}/\texttt{0x6E}),
level to $[\texttt{0x09},\texttt{0x42}]$, the
\texttt{lengthSizeMinusOne}/reserved byte to the form
$\texttt{FC}|x$, and the \texttt{numOfSPS}/reserved byte to
$\texttt{E0}|n$ with $n=1$.  This signature is sufficient to
uniquely identify the record amid the surrounding BIFS commands.  On
the K8B reference capture the recovered record is:

\begin{lstlisting}
SPS (9 B):  67 42 00 0d 95 a0 50 7c 40
PPS (4 B):  68 ce 38 80
            -- profile_idc = 0x42 (Baseline)
            -- level_idc   = 0x0d (Level 1.3)
\end{lstlisting}

Prepending these as Annex-B NAL units to the extracted slice stream
produces a fully self-contained H.264 byte stream that FFmpeg
identifies as \texttt{Baseline 320x240 25\,fps yuv420p}.

% ====================================================================
\section{Reproducibility and Metadata}\label{sec:reproducibility}

\begin{table}[t]
\centering
\caption{Reference capture metadata (excerpt of the JSON sidecar
\texttt{data/captures/k8b\_100pct.json}).}
\label{tab:capture-meta}
\begin{tabular}{ll}
\toprule
Field & Value \\
\midrule
Channel              & K8B \\
Frequency            & 183.008\,MHz \\
Sample rate          & 3\,MSPS \\
Gain mode            & Airspy hardware AGC \\
Bias-tee             & off \\
Antenna              & Brisa flat-panel indoor with USB-powered LNA \\
Antenna spec (claimed)& 30\,dBi composite gain, 5\,m coax feed \\
Location             & Seoul metropolitan area, indoor near window \\
File                 & 30.8\,MB, 5\,014 ETI frames ($\sim$120\,s) \\
FIB success rate     & 100\,\% \\
Ensemble label       & YTN DMB (EId 0xE040) \\
Services decoded     & 5 (mYTN, HD mYTN, 4DRIVE, LOTTE, EWS) \\
\bottomrule
\end{tabular}
\end{table}

Every capture taken with the bundled \texttt{tools/capture\_logged.py}
wrapper produces a JSON sidecar with the schema illustrated in
Table~\ref{tab:capture-meta}.  The reference captures shipped with the
artefact have been retro-fitted with sidecars from build-system records.

% ====================================================================
\section{Related work}\label{sec:related}

The closest predecessor is Eo and Bahk~\cite{eo2024tdmb}, who use
\texttt{gr-dab} for the DAB layer and a patched FFmpeg for the
T-DMB-specific decoding, and who explicitly document several of the
gotchas we re-encountered (deinterleaver sync alignment, missing
composition timestamps, SL packet header parsing).  Their work
remains the canonical reference for the encoding chain.  By contrast
this work emphasises packaging (a single Python package with
documented API) and reproducibility (per-capture metadata sidecars,
a single one-shot decoder script, no run-time FFmpeg modification).

\texttt{welle.io}~\cite{welleio} is a high-quality cross-platform DAB
audio receiver, but does not decode the T-DMB video stream type.
\texttt{eti-cmdline}~\cite{etistuff} (from which we re-use the OFDM
and Viterbi back-end via a small patch) outputs ETI(NI) frames only
and was the starting point for our work.  An earlier portable hardware
implementation is described by Lee~et~al.~\cite{lee2007portable},
and the Korean T-DMB system as a whole is documented in
\cite{ttak0026,etsi102427,etsi102428,hwang2015simulator,iso14496-1}.

% ====================================================================
\section{Conclusion}\label{sec:conclusion}

We have presented, to the best of our knowledge, the first packaged
end-to-end open-source software receiver for Korean Terrestrial DMB.
The pipeline takes a 3\,MSPS I/Q capture from an Airspy~Mini SDR and
produces a playable H.264 video file using only freely-distributable
components.  Two implementation lessons stand out: (i)~the Forney
deinterleaver must be aligned to the TS sync byte \emph{before}
processing rather than after, and (ii)~Korean broadcasters embed the
H.264 AVCC inside the BIFS/SceneDescription stream, not the
ObjectDescriptor stream as MPEG-4 Systems prescribes.  Either lesson
silently breaks an otherwise-correct implementation.  The code,
captures, and metadata schema are released to allow further research
on T-DMB and on its now-encrypted HD~DMB successor.

% ====================================================================
\section*{Code and data availability}
The Python package, capture tooling, reference ETI capture, decoded
IDR thumbnails, and this paper's source live at
\url{https://github.com/zobithecat/airspy-mini-dmb}.  All code is MIT-licensed;
the reference RF capture is released under CC-BY-4.0.

\section*{Acknowledgements}
We rely on, and are grateful for, the prior open-source work of
Jan van Katwijk (\texttt{eti-stuff}), Albrecht Lohofener et al.\
(\texttt{welle.io}), and Sungbo Eo and Saewoong Bahk
(\texttt{dmb-oss}).

\section*{Author contributions and use of AI}
\textbf{Seonggeun Yoo} (primary author) conceived the project,
captured all RF data, made every experimental and design decision,
diagnosed the sync-alignment failure on real captures, validated
results visually, and edited the manuscript.
\textbf{Claude} (Anthropic, large language model) acted as a
paired-programming and drafting collaborator throughout: it wrote
much of the initial Python implementation, reverse-engineered the
Korean T-DMB SL-packet header layout from the raw byte patterns
present in our captures, located the AVC
\texttt{DecoderConfigurationRecord} embedded in the BIFS stream,
and produced the first draft of this paper.  Per the arXiv AI
disclosure guidelines, we emphasise that all code, results, and
prose were reviewed by the primary author before submission and
the primary author takes full responsibility for the work.

% ====================================================================
\printbibliography

\end{document}
